% GPIO chapter. Section order is the per-peripheral template enforced by
% python/check_intro_names.py. Vector numbers are frozen by the generator
% (python/generate.py, _mcuMpIrqVectors).
\section{Overview}

Each \peripheral{GPIOx} peripheral is one 8-pin general-purpose input/output port. A port is a parallel register set, so all eight pins are updated in one access, and every pin is also configurable on its own. Pin $y$ of port $x$ is written P$x$.$y$ and is controlled by bit $y$ of each port register. Every pin is either in GPIO mode, driven by the port registers, or in alternate-function mode, driven by one of up to eight peripheral functions selected per pin.

\begin{itemize}
	\item 8 pins per port, individually configurable as input, output, or resistor-terminated input
	\item Atomic set, clear and toggle registers (\register{PxOUTS}, \register{PxOUTC}, \register{PxOUTT}), so no read-modify-write is needed
	\item Alternate-function mode per pin (\register{PxSEL}) with an 8-way function select per pin (\register{PxAFS})
	\item Relocatable peripheral inputs, routed by \register{PxAFS} independently of \register{PxSEL}
	\item One interrupt per pin, 8 vectors per port, with programmable edge (\register{PxIES})
	\item Event-fabric task select per pin (\register{PxTASK}) in configurations with the event fabric
\end{itemize}

\section{Functional description}

\subsection{GPIO mode}

In GPIO mode the pin is controlled by \register{PxDIR}, \register{PxOUT} and \register{PxREN}: \register{PxDIR} selects input or output, \register{PxOUT} selects the output level, and \register{PxREN} enables the internal pull resistor on an input. Table \ref{t:gpio-config} lists the four states.

\begin{table}[htb]
	\centering
	\caption{GPIO pin states.}
	\label{t:gpio-config}
	\begin{tabular}{ c c c l }
	\textbf{\register{PxDIR}[y]} & \textbf{\register{PxOUT}[y]} & \textbf{\register{PxREN}[y]} & \textbf{P$x$.$y$ state} \\ \hline
	0 & - & 0 & High-impedance input, resistor disabled \\
	\rowcolor{tablehighlightcolor}
	0 & - & 1 & Input with pull resistor enabled \\
	1 & 0 & - & Output driven low \\
	\rowcolor{tablehighlightcolor}
	1 & 1 & - & Output driven high \\
	\hline
	\end{tabular}
\end{table}

\register{PxIN} always reads the digital level of the pad, in both directions and in both modes. The read value is latched on the falling edge of the memory enable signal, so a pin that changes during the access returns one stable sample.

\subsection{Set, clear and toggle registers}

Writing 1 to a bit of \register{PxOUTS}, \register{PxOUTC} or \register{PxOUTT} sets, clears or toggles the corresponding \register{PxOUT} bit; writing 0 has no effect. Each is a single write, so two harts, or a handler and the code it interrupted, can change different pins of one port without a read-modify-write race. \register{PxOUTS} and \register{PxOUTT} read back \register{PxOUT}; \register{PxOUTC} reads back its complement.

\subsection{Alternate-function mode}

Two registers together decide what drives a pin. \register{PxSEL}[y] = 1 puts P$x$.$y$ in alternate-function mode, and the 4-bit field $y$ of \register{PxAFS} (bits $4y+3$ to $4y$, of which the low three are implemented) selects which function, AF0 to AF7, governs it. \register{PxAFS} resets to 0, so a pin in alternate-function mode drives its AF0 function, which is the pin's home peripheral function; software that only writes \register{PxSEL} therefore obtains the home function. Section \ref{s:pinsConfig} tabulates every pin's plane assignment.

In alternate-function mode the selected function owns the pin's direction, output and resistor controls, and \register{PxDIR}[y], \register{PxOUT}[y] and \register{PxREN}[y] have no effect until \register{PxSEL}[y] returns to 0. A field value that selects a plane on which the pin has no function leaves the pin as a high-impedance input. Some peripheral inputs, such as the timer capture pins, are inputs on every plane and are governed by \register{PxDIR} regardless of \register{PxSEL}, because the peripheral never drives them.

The planes above AF0 carry a shared pool of relocatable output functions (the \peripheral{UARTx} transmitters, the \peripheral{TIMERx} compare outputs, the \peripheral{SPI1} clock and master-out) fanned across the ports, so a peripheral output can be placed on whichever pin suits the board. Because there is one instance of each function, selecting it on several pins drives the same signal on all of them. A few plane slots carry bidirectional bus completions so that whole buses relocate as a set: both I2C pairs, the \peripheral{UART0} receiver alongside its transmitter, and the \peripheral{SPI1} master-in alongside its clock and master-out. A function that belongs to a peripheral this configuration omits leaves the pool, and its slots become high-impedance inputs.

\subsection{Relocatable peripheral inputs}

A peripheral input that appears at more than one pin (UART receivers, timer captures, I2C data and clock) is routed from the pin whose \register{PxAFS} field selects its plane, and from its home pin otherwise, regardless of \register{PxSEL}. Input routing ignores \register{PxSEL} so that the routing mirrors \register{PxIN}, which is always visible; a pin can therefore stay in GPIO output mode while its \register{PxAFS} field routes the pad into a peripheral input, which gives a software loopback for self-test. If the same input is selected on more than one pin, a fixed priority (newest alternate location first, then the older alternate, then home) picks one, so the outcome of that configuration error is deterministic.

\subsection{Pin interrupts}

Each pin has its own interrupt. With \register{PxIE}[y] set, the edge selected by \register{PxIES}[y] (0 = rising, 1 = falling) sets \register{PxIF}[y] and raises the pin's vector. Pin interrupts stay available in alternate-function mode alongside any interrupt the governing peripheral raises. \register{PxIF} is latched on the falling edge of the memory enable signal, so a handler reads one stable sample.

\subsection{Event-fabric tasks}

\register{PxTASK} selects which pins respond to the event fabric's OUT-SET and OUT-CLR tasks for this port; in a configuration without the event fabric the register is writable and has no effect. A task pulse is applied after any coincident register write in the same cycle, and CLR after SET, so a same-cycle conflict resolves to the cleared, safe direction.

\section{Interrupts and events}

Port $x$ owns eight consecutive vectors, one per pin, in the order listed in Table \ref{t:gpio-irq}. The interrupt is a level derived from \register{PxIF}[y] and \register{PxIE}[y].

\begin{table}[htb]
	\centering
	\caption{\peripheral{GPIOx} interrupt vectors; pin $y$ uses the port's base vector plus $y$.}
	\label{t:gpio-irq}
	\begin{tabular}{ l l p{3.6cm} p{2.8cm} p{3.0cm} }
	\textbf{Port} & \textbf{Vectors} & \textbf{Condition} & \textbf{Set by} & \textbf{Cleared by} \\ \hline
	\peripheral{GPIO0} & 1 to 8 & Selected edge on P0.$y$ & \register{PxIF}[y] & Writing 1 to \register{PxIF}[y] \\
	\rowcolor{tablehighlightcolor}
	\peripheral{GPIO1} & 28 to 35 & Selected edge on P1.$y$ & \register{PxIF}[y] & Writing 1 to \register{PxIF}[y] \\
	\peripheral{GPIO2} & 36 to 43 & Selected edge on P2.$y$ & \register{PxIF}[y] & Writing 1 to \register{PxIF}[y] \\
	\rowcolor{tablehighlightcolor}
	\peripheral{GPIO3} & 44 to 51 & Selected edge on P3.$y$ & \register{PxIF}[y] & Writing 1 to \register{PxIF}[y] \\
	\peripheral{GPIO4} & 98 to 105 & Selected edge on P4.$y$ & \register{PxIF}[y] & Writing 1 to \register{PxIF}[y] \\
	\rowcolor{tablehighlightcolor}
	\peripheral{GPIO5} & 106 to 113 & Selected edge on P5.$y$ & \register{PxIF}[y] & Writing 1 to \register{PxIF}[y] \\
	\hline
	\end{tabular}
\end{table}

\section{Programming}

\subsection{Configuring a pin as an output}

\begin{enumerate}
	\item Write the initial level to \register{PxOUT}[y].
	\item Clear \register{PxSEL}[y].
	\item Set \register{PxDIR}[y].
\end{enumerate}

\subsection{Routing a peripheral output to an alternate pin}

\begin{enumerate}
	\item Write the function's plane number to field $y$ of \register{PxAFS}.
	\item Set \register{PxSEL}[y].
\end{enumerate}

\subsection{Enabling a pin interrupt}

\begin{enumerate}
	\item Clear \register{PxIE}[y].
	\item Write the edge to \register{PxIES}[y].
	\item Write 1 to \register{PxIF}[y] to discard any stale flag.
	\item Enable the pin's vector in the hart's \peripheral{IRQROUTER} row.
	\item Set \register{PxIE}[y].
\end{enumerate}

\section{Usage notes}

\paragraph{Disable the pin interrupt before changing its edge.} Writing \register{PxIES}[y] while \register{PxIE}[y] is set can raise a spurious interrupt, because the edge detector sees the selection change as an edge.

\paragraph{Clear \register{PxIF}[y] with a 1 before returning.} The interrupt is a level; a handler that returns with the flag set is re-entered.

\paragraph{Configure capture pins as inputs in \register{PxDIR}.} A timer capture input is governed by \register{PxDIR} on every plane, so setting \register{PxSEL}[y] alone does not make the pad an input.

\paragraph{Select an input function on one pin only.} The priority among several selected pins is fixed and makes the error deterministic, not correct.

\paragraph{Use the set, clear and toggle registers from shared code.} A read-modify-write of \register{PxOUT} from two harts loses one of the two updates.
